\section*{الفصل \ref{ch:g3:continents-and-oceans} --- مجسم الكرة الأرضية: القارات والمحيطات}

\begin{solution}{exo:g3:continents-and-oceans:1}
آسيا، وإفريقيا، وأمريكا الشمالية، وأمريكا الجنوبية، والقارة القطبية الجنوبية، وأوروبا،
وأستراليا. أكبرها آسيا؛ وأصغرها أستراليا.
\end{solution}

\begin{solution}{exo:g3:continents-and-oceans:2}
\omterm{def:g2:water-on-earth:ocean}{المحيط} الهادئ، والأطلسي، والهندي، والجنوبي، والمتجمد الشمالي.
أكبرها \omterm{def:g2:water-on-earth:ocean}{المحيط} الهادئ؛ وأصغرها \omterm{def:g2:water-on-earth:ocean}{المحيط} المتجمد الشمالي.
\end{solution}

\begin{solution}{exo:g3:continents-and-oceans:3}
\omterm{def:g3:continents-and-oceans:poles}{القطب الشمالي} في أعلى \omterm{def:g3:continents-and-oceans:globe}{مجسم الكرة الأرضية}؛ \omterm{def:g3:continents-and-oceans:poles}{والقطب الجنوبي} في أسفله؛
\omterm{def:g3:continents-and-oceans:equator}{وخط الاستواء} يدور حول وسطه كله.
\end{solution}

\begin{solution}{exo:g3:continents-and-oceans:4}
أوروبا: الشمالي. أستراليا: الجنوبي. القارة القطبية الجنوبية: الجنوبي. أمريكا الشمالية:
الشمالي.
\end{solution}

\begin{solution}{exo:g3:continents-and-oceans:5}
\omterm{def:g2:water-on-earth:ocean}{المحيط} المتجمد الشمالي يقع حول \omterm{def:g3:continents-and-oceans:poles}{القطب الشمالي}. \omterm{def:g2:water-on-earth:ocean}{والمحيط} الجنوبي يقع حول
القارة القطبية الجنوبية.
\end{solution}

\begin{solution}{exo:g3:continents-and-oceans:6}
من إفريقيا إلى أمريكا الجنوبية: \omterm{def:g2:water-on-earth:ocean}{المحيط} الأطلسي. ومن أستراليا شرقًا إلى
أمريكا الجنوبية: \omterm{def:g2:water-on-earth:ocean}{المحيط} الهادئ. ومن إفريقيا إلى أستراليا: \omterm{def:g2:water-on-earth:ocean}{المحيط}
الهندي.
\end{solution}

\begin{solution}{exo:g3:continents-and-oceans:7}
يعيش سام قرب \omterm{def:g3:continents-and-oceans:equator}{خط الاستواء}، حيث الحر طوال السنة. ويعيش بو
قرب \omterm{def:g3:continents-and-oceans:poles}{قطب}، حيث البرد طوال السنة.
\end{solution}

\begin{solution}{exo:g3:continents-and-oceans:8}
$10 - 7 = 3$: يقع إصبعها على اليابسة 3 مرات. وهذا ما نستطيع أن
نتوقعه: نحو 71 مربعًا من كل 100 من الأرض ماء، إذن يقف الإصبع على \omterm{def:g2:water-on-earth:ocean}{البحر} نحو 7 مرات
من كل 10. (وبعشر دورات فقط، قد تحصل أيضًا على 6 أو 8: فالعدد يتغير من لعبة
إلى أخرى.)
\end{solution}

\begin{solution}{exo:g3:continents-and-oceans:9}
لرسم \omterm{def:g3:continents-and-oceans:world-map}{خريطة العالم}، قُطع \omterm{def:g3:continents-and-oceans:globe}{المجسم} على طول خط يمر بوسط \omterm{def:g2:water-on-earth:ocean}{المحيط} الهادئ ثم
بُسط. فانتهى نصف \omterm{def:g2:water-on-earth:ocean}{المحيط} الهادئ على الحافة اليسرى \omterm{def:g2:drawing-to-map:map}{للخريطة}، ونصفه الآخر
على الحافة اليمنى. وليس على الأرض الحقيقية إلا \omterm{def:g2:water-on-earth:ocean}{محيط} هادئ واحد: فحافتا
\omterm{def:g2:drawing-to-map:map}{الخريطة} هما الخط نفسه على \omterm{def:g3:continents-and-oceans:globe}{المجسم}.
\end{solution}

\begin{solution}{exo:g3:continents-and-oceans:10}
توم مخطئ. \omterm{def:g3:continents-and-oceans:continent}{فالقارات} من الأكبر إلى الأصغر هي آسيا، وإفريقيا، وأمريكا
الشمالية، وأمريكا الجنوبية، والقارة القطبية الجنوبية، وأوروبا، وأستراليا:
فالقارة القطبية الجنوبية الخامسة فقط، وأصغر بكثير من آسيا. وتخدعه \omterm{def:g3:continents-and-oceans:world-map}{خريطة
العالم} لأن \omterm{def:g2:drawing-to-map:map}{الخريطة} المسطحة تمدد اليابسة قرب \omterm{def:g3:continents-and-oceans:poles}{القطبين} كثيرًا.
فالقارة القطبية الجنوبية، الواقعة حول \omterm{def:g3:continents-and-oceans:poles}{القطب الجنوبي}، تمددت فصارت شريطًا
طويلًا على طول أسفل \omterm{def:g2:drawing-to-map:map}{الخريطة} كلها. أما على \omterm{def:g3:continents-and-oceans:globe}{المجسم} فهي كتلة مستديرة من
اليابسة، أصغر بكثير من آسيا.
\end{solution}
